% GENERATED FILE. Edit canonical lessons, then run npm run generate:books.
% canonical-source-hash: fnv1a64:bd7a88489c00989f
% canonical-lessons: TE-C525-upadhi, TE-C525-sampadana, TE-C525-adayam, TE-C525-vetanam, TE-C525-sramincu, TE-R525-review-518-525

\chapter{Making a Living}
\label{ch:te-making-a-living}
\hlchaptermodality{525}

\begin{chapteropening}
\textbf{By the end of this chapter:} \emph{I can talk about livelihood, earnings, income, wages and hard work.}

Complete the last lesson of chapter 525: I can talk about livelihood, earnings, income, wages and hard work.
\end{chapteropening}

\section[\texorpdfstring{\te{ఉపాధి} (upādhi)}{ఉపాధి (upādhi)}]{\te{ఉపాధి} (upādhi) — a livelihood; employment}

\label{lesson:TE-C525-upadhi}



\begin{quote}
Before the new one: say the Telugu for a skilled person, then the Telugu for initiative.
\end{quote}

\subsection*{You'll want to know: \te{ఉపాధి}}
\textbf{\te{ఉపాధి}} — \emph{upādhi} — \textquotedblleft{}a livelihood; employment\textquotedblright{}.

\begin{quote}
\te{పల్లెల్లో} \te{యువకులకు} \te{ఉపాధి} \te{దొరకడం} \te{కష్టం}.
\end{quote}

\begin{quote}
\emph{pallellō yuvakulaku upādhi dorakaḍaṁ kaṣṭaṁ.}
\end{quote}

\begin{quote}
In the villages it is hard for young people to find work.
\end{quote}

\te{ఉపాధి} \te{అవకాశాలు}, job opportunities.

\begin{cousinweb}
\textbf{\te{ఉపాధి}} (\emph{upādhi}) is borrowed from Sanskrit \emph{upādhi} (CDIAL 2299), \textquotedblleft{}a point of view, an aim\textquotedblright{}; in Telugu it means a means of living.
\end{cousinweb}

\begin{grammarlens}[title={What you've built}]
One more word for describing what you lived and felt.
\end{grammarlens}

\subsection*{Your turn}
\begin{itemize}
\raggedright
  \item \emph{Say it:} \emph{upādhi}
  \item \emph{Say it:} \emph{upādhi}, once more
  \item \emph{Say it:} \emph{pallellō yuvakulaku upādhi dorakaḍaṁ kaṣṭaṁ.} — In the villages it is hard for young people to find work.
  \item \emph{Recall:} say the Telugu for a skilled person, then the Telugu for initiative, then say \emph{upādhi} again
\end{itemize}

\begin{tcolorbox}[breakable,colback=teal!4,colframe=teal!35!black,title={Before you move on}]
What does \te{ఉపాధి} mean? (\textquotedblleft{}A livelihood; employment\textquotedblright{}.) Say it once more.
\end{tcolorbox}
\section[\texorpdfstring{\te{సంపాదన} (sampādana)}{సంపాదన (sampādana)}]{\te{సంపాదన} (sampādana) — earnings}

\label{lesson:TE-C525-sampadana}



\begin{quote}
Before the new one: say the Telugu for initiative, then the Telugu for a livelihood.
\end{quote}

\subsection*{You'll want to know: \te{సంపాదన}}
\textbf{\te{సంపాదన}} — \emph{sampādana} — \textquotedblleft{}earnings\textquotedblright{}.

\begin{quote}
\te{అతని} \te{సంపాదనతో} \te{కుటుంబం} \te{గడుస్తుంది}.
\end{quote}

\begin{quote}
\emph{atani sampādanatō kuṭumbaṁ gaḍustundi.}
\end{quote}

\begin{quote}
The family lives on his earnings.
\end{quote}

The money a person earns; its verb is \te{సంపాదించు}, to earn.

\begin{cousinweb}
\textbf{\te{సంపాదన}} (\emph{sampādana}) is borrowed from Sanskrit \emph{sampādana}, \textquotedblleft{}procuring; accomplishing\textquotedblright{} (Monier-Williams).
\end{cousinweb}

\begin{grammarlens}[title={What you've built}]
One more word for describing what you lived and felt.
\end{grammarlens}

\subsection*{Your turn}
\begin{itemize}
\raggedright
  \item \emph{Say it:} \emph{sampādana}
  \item \emph{Say it:} \emph{sampādana}, once more
  \item \emph{Say it:} \emph{atani sampādanatō kuṭumbaṁ gaḍustundi.} — The family lives on his earnings.
  \item \emph{Recall:} say the Telugu for initiative, then the Telugu for a livelihood, then say \emph{sampādana} again
\end{itemize}

\begin{tcolorbox}[breakable,colback=teal!4,colframe=teal!35!black,title={Before you move on}]
What does \te{సంపాదన} mean? (\textquotedblleft{}Earnings\textquotedblright{}.) Say it once more.
\end{tcolorbox}
\section[\texorpdfstring{\te{ఆదాయం} (ādāyaṁ)}{ఆదాయం (ādāyaṁ)}]{\te{ఆదాయం} (ādāyaṁ) — income}

\label{lesson:TE-C525-adayam}



\begin{quote}
Before the new one: say the Telugu for a livelihood, then the Telugu for earnings.
\end{quote}

\subsection*{You'll want to know: \te{ఆదాయం}}
\textbf{\te{ఆదాయం}} — \emph{ādāyaṁ} — \textquotedblleft{}income\textquotedblright{}.

\begin{quote}
\te{పంట} \te{బాగా} \te{పండితే} \te{రైతుల} \te{ఆదాయం} \te{పెరుగుతుంది}.
\end{quote}

\begin{quote}
\emph{paṇṭa bāgā paṇḍitē raitula ādāyaṁ perugutundi.}
\end{quote}

\begin{quote}
If the harvest is good, the farmers' income rises.
\end{quote}

The opposite of \te{ఖర్చు}, expense.

\begin{cousinweb}
\textbf{\te{ఆదాయం}} (\emph{ādāyaṁ}) is usually traced to Sanskrit \emph{ādāya}, \textquotedblleft{}taking, seizing\textquotedblright{} (Monier-Williams); compare Sanskrit \emph{āya}, \textquotedblleft{}arrival; income\textquotedblright{} (CDIAL 1283).
\end{cousinweb}

\begin{grammarlens}[title={What you've built}]
One more word for describing what you lived and felt.
\end{grammarlens}

\subsection*{Your turn}
\begin{itemize}
\raggedright
  \item \emph{Say it:} \emph{ādāyaṁ}
  \item \emph{Say it:} \emph{ādāyaṁ}, once more
  \item \emph{Say it:} \emph{paṇṭa bāgā paṇḍitē raitula ādāyaṁ perugutundi.} — If the harvest is good, the farmers' income rises.
  \item \emph{Recall:} say the Telugu for a livelihood, then the Telugu for earnings, then say \emph{ādāyaṁ} again
\end{itemize}

\begin{tcolorbox}[breakable,colback=teal!4,colframe=teal!35!black,title={Before you move on}]
What does \te{ఆదాయం} mean? (\textquotedblleft{}Income\textquotedblright{}.) Say it once more.
\end{tcolorbox}
\section[\texorpdfstring{\te{వేతనం} (vētanaṁ)}{వేతనం (vētanaṁ)}]{\te{వేతనం} (vētanaṁ) — a wage; pay}

\label{lesson:TE-C525-vetanam}



\begin{quote}
Before the new one: say the Telugu for earnings, then the Telugu for income.
\end{quote}

\subsection*{You'll want to know: \te{వేతనం}}
\textbf{\te{వేతనం}} — \emph{vētanaṁ} — \textquotedblleft{}a wage; pay\textquotedblright{}.

\begin{quote}
\te{కార్మికులు} \te{ఎక్కువ} \te{వేతనం} \te{కోరారు}.
\end{quote}

\begin{quote}
\emph{kārmikulu ekkuva vētanaṁ kōrāru.}
\end{quote}

\begin{quote}
The workers asked for higher wages.
\end{quote}

Formal; \te{జీతం} is the everyday word.

\begin{cousinweb}
\textbf{\te{వేతనం}} (\emph{vētanaṁ}) is borrowed from Sanskrit \emph{vetana}, \textquotedblleft{}wages, hire, salary, livelihood\textquotedblright{} (Monier-Williams).
\end{cousinweb}

\begin{grammarlens}[title={What you've built}]
One more word for describing what you lived and felt.
\end{grammarlens}

\subsection*{Your turn}
\begin{itemize}
\raggedright
  \item \emph{Say it:} \emph{vētanaṁ}
  \item \emph{Say it:} \emph{vētanaṁ}, once more
  \item \emph{Say it:} \emph{kārmikulu ekkuva vētanaṁ kōrāru.} — The workers asked for higher wages.
  \item \emph{Recall:} say the Telugu for earnings, then the Telugu for income, then say \emph{vētanaṁ} again
\end{itemize}

\begin{tcolorbox}[breakable,colback=teal!4,colframe=teal!35!black,title={Before you move on}]
What does \te{వేతనం} mean? (\textquotedblleft{}A wage; pay\textquotedblright{}.) Say it once more.
\end{tcolorbox}
\section[\texorpdfstring{\te{శ్రమించు} (śramiñcu)}{శ్రమించు (śramiñcu)}]{\te{శ్రమించు} (śramiñcu) — to toil; to labour}

\label{lesson:TE-C525-sramincu}



\begin{quote}
Before the new one: say the Telugu for income, then the Telugu for a wage.
\end{quote}

\subsection*{You'll want to know: \te{శ్రమించు}}
\textbf{\te{శ్రమించు}} — \emph{śramiñcu} — \textquotedblleft{}to toil; to labour\textquotedblright{}.

\begin{quote}
\te{ఆ} \te{పుస్తకం} \te{కోసం} \te{రచయిత} \te{రెండేళ్ళు} \te{శ్రమించారు}.
\end{quote}

\begin{quote}
\emph{ā pustakaṁ kōsaṁ racayita reṇḍēḷḷu śramiñcāru.}
\end{quote}

\begin{quote}
The writer toiled for two years on that book.
\end{quote}

Past (I …): \te{శ్రమించాను}, \emph{śramiñcānu}.

The noun is \te{శ్రమ}, toil.

\begin{cousinweb}
\textbf{\te{శ్రమించు}} (\emph{śramiñcu}) is built on Sanskrit \emph{śrama}, \textquotedblleft{}labour; fatigue\textquotedblright{} (CDIAL 12682), with the Telugu verb ending -\te{ఇంచు}.
\end{cousinweb}

\begin{grammarlens}[title={What you've built}]
One more word for describing what you lived and felt.
\end{grammarlens}

\subsection*{Your turn}
\begin{itemize}
\raggedright
  \item \emph{Say it:} \emph{śramiñcu}
  \item \emph{Say it:} \emph{śramiñcu}, once more
  \item \emph{Say it:} \emph{ā pustakaṁ kōsaṁ racayita reṇḍēḷḷu śramiñcāru.} — The writer toiled for two years on that book.
  \item \emph{Recall:} say the Telugu for income, then the Telugu for a wage, then say \emph{śramiñcu} again
\end{itemize}

\begin{tcolorbox}[breakable,colback=teal!4,colframe=teal!35!black,title={Before you move on}]
What does \te{శ్రమించు} mean? (\textquotedblleft{}To toil; to labour\textquotedblright{}.) Say it once more.
\end{tcolorbox}
\section[(dialogue)]{Words from chapters 518–525 — a review}

\label{lesson:TE-R525-review-518-525}



\begin{quote}
Say the Telugu for to get hurt and a bump.
\end{quote}

\subsection*{Your turn}
Say these aloud:

\begin{itemize}
\raggedright
  \item to get hurt, a bump, to stagger, to limp, bleeding
  \item treatment, surgery, an anaesthetic, a bandage, to visit the sick
  \item to steady, to grow strong, to be worn out, to blow one's nose, to pick up
  \item frailty, to grow weak, to grow thin, to put on weight, nutrition
  \item a profession, status, a position, a promotion, an ambition
  \item a worker, an expert, an adviser, an assistant, a contractor
  \item talent, proficiency, workmanship, a skilled person, initiative
  \item a livelihood, earnings, income, a wage, to toil
\end{itemize}

\begin{tcolorbox}[breakable,colback=teal!4,colframe=teal!35!black,title={Before you move on}]
To get hurt? (\textbf{\te{గాయపడు}}, \emph{gāyapaḍu}.) A bump? (\textbf{\te{బొప్పి}}, \emph{boppi}.) To stagger? (\textbf{\te{తూలు}}, \emph{tūlu}.) To limp? (\textbf{\te{కుంటు}}, \emph{kuṇṭu}.) Bleeding? (\textbf{\te{రక్తస్రావం}}, \emph{raktasrāvaṁ}.) Treatment? (\textbf{\te{చికిత్స}}, \emph{cikitsa}.) Surgery? (\textbf{\te{శస్త్రచికిత్స}}, \emph{śastracikitsa}.) An anaesthetic? (\textbf{\te{మత్తుమందు}}, \emph{mattumandu}.) A bandage? (\textbf{\te{పట్టీ}}, \emph{paṭṭī}.) To visit the sick? (\textbf{\te{పరామర్శించు}}, \emph{parāmarśiñcu}.) To steady? (\textbf{\te{కుదుటపడు}}, \emph{kuduṭapaḍu}.) To grow strong? (\textbf{\te{బలపడు}}, \emph{balapaḍu}.) To be worn out? (\textbf{\te{నలిగిపోవు}}, \emph{naligipōvu}.) To blow one's nose? (\textbf{\te{చీదు}}, \emph{cīdu}.) To pick up? (\textbf{\te{పుంజుకో}}, \emph{puñjukō}.) Frailty? (\textbf{\te{బలహీనత}}, \emph{balahīnata}.) To grow weak? (\textbf{\te{నీరసించు}}, \emph{nīrasiñcu}.) To grow thin? (\textbf{\te{చిక్కిపోవు}}, \emph{cikkipōvu}.) To put on weight? (\textbf{\te{లావెక్కు}}, \emph{lāvekku}.) Nutrition? (\textbf{\te{పోషకాహారం}}, \emph{pōṣakāhāraṁ}.) A profession? (\textbf{\te{వృత్తి}}, \emph{vṛtti}.) Status? (\textbf{\te{హోదా}}, \emph{hōdā}.) A position? (\textbf{\te{పదవి}}, \emph{padavi}.) A promotion? (\textbf{\te{పదోన్నతి}}, \emph{padōnnati}.) An ambition? (\textbf{\te{ఆశయం}}, \emph{āśayaṁ}.) A worker? (\textbf{\te{కార్మికుడు}}, \emph{kārmikuḍu}.) An expert? (\textbf{\te{నిపుణుడు}}, \emph{nipuṇuḍu}.) An adviser? (\textbf{\te{సలహాదారు}}, \emph{salahādāru}.) An assistant? (\textbf{\te{సహాయకుడు}}, \emph{sahāyakuḍu}.) A contractor? (\textbf{\te{గుత్తేదారు}}, \emph{guttēdāru}.) Talent? (\textbf{\te{ప్రతిభ}}, \emph{pratibha}.) Proficiency? (\textbf{\te{ప్రావీణ్యం}}, \emph{prāvīṇyaṁ}.) Workmanship? (\textbf{\te{పనితనం}}, \emph{panitanaṁ}.) A skilled person? (\textbf{\te{నేర్పరి}}, \emph{nērpari}.) Initiative? (\textbf{\te{చొరవ}}, \emph{corava}.) A livelihood? (\textbf{\te{ఉపాధి}}, \emph{upādhi}.) Earnings? (\textbf{\te{సంపాదన}}, \emph{sampādana}.) Income? (\textbf{\te{ఆదాయం}}, \emph{ādāyaṁ}.) A wage? (\textbf{\te{వేతనం}}, \emph{vētanaṁ}.) To toil? (\textbf{\te{శ్రమించు}}, \emph{śramiñcu}.)
\end{tcolorbox}
